\documentclass[12pt]{article}
\usepackage[utf8]{inputenc}
\usepackage{amsmath}
\usepackage{amssymb}
\usepackage{graphicx}
\usepackage{booktabs}
\usepackage{listings}
\usepackage{xcolor}
\usepackage{hyperref}
\usepackage{geometry}
\geometry{margin=1in}

\lstset{
    language=Python,
    basicstyle=\ttfamily\small,
    keywordstyle=\color{blue}\bfseries,
    commentstyle=\color{green!60!black},
    stringstyle=\color{red},
    numbers=left,
    numberstyle=\tiny\color{gray},
    stepnumber=1,
    numbersep=5pt,
    backgroundcolor=\color{gray!10},
    frame=single,
    breaklines=true,
    breakatwhitespace=false,
    tabsize=4,
    showstringspaces=false
}

\title{Smart Loaders Documentation}
\author{}
\date{}

\begin{document}

\maketitle

\section{Introduction}

The \texttt{smart\_loaders.py} module provides input abstraction for CTCC: it chooses \texttt{yaml\_loaders} or \texttt{json\_loaders} at import time based on the \texttt{CTCC\_INPUT\_MODE} environment variable, then re-exports all \texttt{load\_*} and related symbols from the chosen backend. Calculation scripts import \texttt{from smart\_loaders import load\_project\_technical\_details}, \texttt{load\_financing\_details}, etc., and get the same API regardless of whether input is YAML or JSON; no code changes are needed when switching mode.

\section{Method Overview}

\subsection{Purpose}

\begin{enumerate}
    \item At import time: read \texttt{CTCC\_INPUT\_MODE} (default \texttt{yaml}); if \texttt{json}, \texttt{from json\_loaders import *}; else \texttt{from yaml\_loaders import *}.
    \item Expose \texttt{get\_input\_mode()} so callers can query the current mode (\texttt{"json"} or \texttt{"yaml"}).
    \item Provide mode-aware raw-data helpers: \texttt{get\_project\_data\_raw()}, \texttt{get\_financing\_data\_raw()}, \texttt{get\_physical\_data\_raw()} (return the raw dict for project technical details, financing, or physical details from either YAML or JSON).
    \item Provide \texttt{load\_terrain\_miles()} that uses \texttt{get\_physical\_data\_raw()} and returns \texttt{terrain["terrain\_miles"]}.
    \item Re-export all \texttt{load\_*} and other symbols from the active backend so scripts use a single import.
\end{enumerate}

\subsection{Calculation Approach}

Dispatch only; no cost/benefit formulas. The module does not implement load logic; it only selects the backend and re-exports. All \texttt{load\_*} implementations live in \texttt{yaml\_loaders.py} or \texttt{json\_loaders.py} (see their documentation).

\subsection{Inputs and Outputs}

\textbf{Inputs:}
\begin{itemize}
    \item \texttt{CTCC\_INPUT\_MODE} environment variable (read once at import): \texttt{yaml} (default) or \texttt{json}.
\end{itemize}

\textbf{Outputs:}
\begin{itemize}
    \item Same \texttt{load\_*} API and return types as \texttt{yaml\_loaders} or \texttt{json\_loaders}; callers receive dataclasses or dicts from the chosen backend.
\end{itemize}

\subsection{Integration with Other Scripts}

\begin{itemize}
    \item Every calculation script that loads configuration imports from \texttt{smart\_loaders} (e.g.\ \texttt{from smart\_loaders import load\_project\_technical\_details}, \texttt{load\_financing\_details}).
    \item \texttt{smart\_loaders} imports either \texttt{json\_loaders} or \texttt{yaml\_loaders} at import time; it does not call both in the same run.
\end{itemize}

\section{Key Functions}

\subsection{Module-level logic (conditional import)}

At import time:
\begin{lstlisting}
_input_mode = os.environ.get('CTCC_INPUT_MODE', 'yaml').lower()
if _input_mode == 'json':
    from json_loaders import *
else:
    from yaml_loaders import *
\end{lstlisting}
All \texttt{load\_*} and other public symbols from the selected module are then available as \texttt{smart\_loaders.load\_project\_technical\_details}, etc.

\subsection{\texttt{get\_input\_mode()}}

Returns the current input mode: \texttt{"json"} or \texttt{"yaml"} (the value determined at import time).

\subsection{\texttt{get\_project\_data\_raw()}}

Returns the raw dict for project technical details (keys \texttt{project}, \texttt{timeline}). In JSON mode calls \texttt{json\_loaders.\_data\_source.get\_data("01\_project\_technical\_details")}; in YAML mode loads \texttt{YAMLS\_DIR / 01\_project\_technical\_details.yaml}.

\subsection{\texttt{get\_financing\_data\_raw()}}

Returns the raw financing dict. In JSON mode \texttt{get\_data("03\_financing")}; in YAML mode loads \texttt{03\_financing.yaml}.

\subsection{\texttt{get\_physical\_data\_raw()}}

Returns the raw physical details dict (includes \texttt{terrain} with \texttt{terrain\_miles}, etc.). In JSON mode \texttt{get\_data("02\_project\_physical\_details")}; in YAML mode loads \texttt{02\_project\_physical\_details.yaml}.

\subsection{\texttt{load\_terrain\_miles()}}

Calls \texttt{get\_physical\_data\_raw()}, then returns \texttt{physical\_details["terrain"]["terrain\_miles"]} (dict mapping terrain type to miles). Raises \texttt{KeyError} if \texttt{terrain} or \texttt{terrain\_miles} is missing.

\subsection{Re-exported \texttt{load\_*} (see yaml/json loader docs)}

All other \texttt{load\_*} functions (e.g.\ \texttt{load\_project\_technical\_details}, \texttt{load\_financing\_details}, \texttt{load\_congestion\_curtailment\_reductions}, \texttt{load\_physical\_details}, \texttt{load\_primary\_bcr\_config}) are defined in \texttt{yaml\_loaders.py} or \texttt{json\_loaders.py} and re-exported here; their parameters, return values, and key logic are documented in those modules.

\section{Example}

Script uses the same imports regardless of mode:
\begin{lstlisting}
from smart_loaders import load_project_technical_details, load_financing_details
details = load_project_technical_details()
financing = load_financing_details()
\end{lstlisting}

To force JSON mode (e.g. before importing):
\begin{lstlisting}
import os
os.environ["CTCC_INPUT_MODE"] = "json"
# Then import smart_loaders (or run script with CTCC_INPUT_MODE=json set in shell)
from smart_loaders import load_project_technical_details
details = load_project_technical_details()  # data from JSON
\end{lstlisting}

\end{document}
